% --- Free Response 1: Market Power & Price Discrimination ---
\question{Starlight Cinemas is the only movie theater in town. They face an inverse demand curve for movie tickets given by \textbf{$P = [[AGG_EQ]]$}, which implies a marginal revenue curve of \textbf{$MR = [[AGG_MR]]$}. The theater’s marginal cost for each additional ticket sold is constant at \textbf{$MC = [[MC_VAL]]$}. Assume the theater has no fixed costs.}

\begin{subquestions}
  \item Assuming Starlight Cinemas charges a single price to all customers, what is the profit-maximizing quantity ($Q^*$) and price ($P^*$)? Show your work and box your final answer.

  \item Calculate the deadweight loss (DWL) in the single-price scenario. Show your work and box your final answer.

  \item Calculate the theater’s total profit at the single price. Show work; box answer.

  \item Now suppose the theater’s demand is actually made up of two distinct groups:

  \begin{itemize}
    \item Locals, whose demand is given by \textbf{$P_L = [[LOC_EQ]]$} (implying $MR_L = [[LOC_MR]]$), and
    \item Visitors, whose demand is given by \textbf{$P_V = [[VIS_EQ]]$} (implying $MR_V = [[VIS_MR]]$).
    \item \textit{Note: Adding these two demands together yields the exact ``aggregate'' demand given initially.}
  \end{itemize}

  If the firm \textit{cannot} distinguish between Locals and Visitors, what price should it charge? How does it compare to the single-price implied by the above parts? Explain briefly. Box your answer.

  \item If the firm \textit{can} successfully distinguish between the two types of customers (e.g., by checking for a local zip code on their ID) and prevent resale, what price will they charge each group, and what is the new total profit they could achieve? Show work; box answer.

  \item If the firm can distinguish between the two types of customers, does the \textit{efficient price} (i.e., the price that would lead to an efficient outcome) change? Briefly explain why or why not.
\end{subquestions}